
% =====================================================================
\section{Laboratorio}
% =====================================================================

\begin{frame}[fragile]{Configuración del entorno}
\begin{lstlisting}[language=bash]
# Opcion A: conda
conda create -n ans python=3.12
conda activate ans
pip install numpy scipy pandas scikit-learn matplotlib seaborn jupyterlab umap-learn pyod

# Opcion B: uv
uv venv && source .venv/bin/activate
uv pip install numpy scipy pandas scikit-learn matplotlib seaborn jupyterlab umap-learn pyod

# Verificar
python -c "import sklearn; print(sklearn.__version__)"
jupyter lab
\end{lstlisting}
  \small Alternativa sin instalación: Google Colab. Crear un repositorio en GitHub para los notebooks del curso.
\end{frame}

\begin{frame}[fragile]{Exploración de un conjunto de datos sin etiquetas}
\begin{lstlisting}
import pandas as pd
from sklearn.datasets import load_wine
from sklearn.pipeline import make_pipeline
from sklearn.impute import SimpleImputer
from sklearn.preprocessing import StandardScaler
from sklearn.metrics import pairwise_distances

df = load_wine(as_frame=True).data           # ignoramos la etiqueta a proposito
print(df.shape); print(df.describe().T)      # escalas muy distintas
print(df.isna().mean())                      # proporcion de faltantes

prep = make_pipeline(SimpleImputer(strategy="median"), StandardScaler())
X = prep.fit_transform(df)

for metrica in ["euclidean", "manhattan", "cosine"]:
    D = pairwise_distances(X, metric=metrica)
    print(metrica, D.mean().round(2))

df.corr().style.background_gradient(cmap="coolwarm")  # correlaciones
\end{lstlisting}
\end{frame}

\begin{frame}{Actividad y tarea}
  \begin{block}{En clase}
    \begin{enumerate}
      \item Instalar el entorno y crear el repositorio de notebooks.
      \item Explorar el conjunto de vinos: escalas, faltantes, correlaciones, histogramas.
      \item Comparar los 5 vecinos más cercanos de una observación con y sin estandarizar, usando distancia euclidiana y coseno.
    \end{enumerate}
  \end{block}
  \begin{exampleblock}{Para la sesión 2}
    \begin{itemize}
      \item Subir el notebook de la sesión al repositorio.
      \item Leer: James et al.\ (2023), cap.~12.1 y 12.4.1 (k-medias); Bishop (2006), cap.~9.1.
      \item Repasar distancia euclidiana y estandarización: k-medias depende de ambas.
    \end{itemize}
  \end{exampleblock}
\end{frame}
